\documentclass[11pt]{article}

\usepackage[margin=1in]{geometry}
\usepackage{amsmath}
\usepackage{booktabs}
\usepackage{hyperref}
\usepackage{listings}
\usepackage{xcolor}

\lstset{
  basicstyle=\ttfamily\small,
  breaklines=true,
  frame=single,
  columns=fullflexible
}

\title{Planck Moment Network: Current Design and Implementation}
\author{}
\date{\today}

\begin{document}

\maketitle

\section{Final Goal}

The final goal is to train one shared neural network using the
component-separated Planck patches produced by
\texttt{Planck Project/compsep\_planck.py}.

Each Planck patch is treated as a realization drawn from a common underlying
distribution.  For every patch, the component-separated Q and U maps are used
as reference clean maps.  Additional clean maps are synthesized from their
scattering statistics, contaminated with nuisance realizations belonging to
the same patch, and used to form supervised training pairs:

\begin{equation}
  \mathrm{input} = \mathrm{synthetic\ clean\ map} + \mathrm{same\mbox{-}patch\ nuisance},
\end{equation}

\begin{equation}
  \mathrm{target} = \mathrm{synthetic\ clean\ map}.
\end{equation}

The training pairs from all patches will eventually be pooled to train a
single shared moment network.  At inference time, this network should produce
posterior moment maps, particularly posterior mean maps, for every patch.
These final maps must have the same spatial dimensions as the original Planck
patches.

\section{Relevant Existing Code}

The implementation was developed after examining:

\begin{itemize}
  \item \texttt{STL\_main/}, especially \texttt{Synthesis.py},
        \texttt{ST\_Operator.py}, and the FFT and kernel data classes;
  \item \texttt{Planck Project/compsep\_planck.py};
  \item \texttt{Moment Networks Main/};
  \item \texttt{Moment Networks Pierre/}; and
  \item \texttt{moment\_network\_notes.txt}.
\end{itemize}

The current \texttt{STL\_main/Synthesis.py} exposes
\texttt{optimize\_from\_maps} and \texttt{synthesize\_from\_maps}.  The older
Moment Network scripts import an older symbol named
\texttt{optimize\_scattering\_LBFGS}, which is not present in the current
checkout.  The new implementation therefore uses
\texttt{optimize\_from\_maps}.

\section{Current Files}

The current implementation is located in:

\begin{lstlisting}
Planck Project/Moment Network/
    make_dataset.py
    make_dataset.sbatch
\end{lstlisting}

\texttt{make\_dataset.py} performs synthesis, augmentation, nuisance addition,
and per-patch dataset-shard creation.  It does not yet train the final shared
moment network.

\texttt{make\_dataset.sbatch} launches the dataset generation on a NERSC GPU
node.  Its parameters are written explicitly in the \texttt{export} section so
that a run's configuration can be read and changed directly in the file.

\section{Input Discovery}

Component-separation results are discovered from:

\begin{lstlisting}
/pscratch/sd/a/atsouros/STL/planck_results/version_2
\end{lstlisting}

The directory contains one NumPy result and one JSON metadata file per patch,
with names such as:

\begin{lstlisting}
p0_m384_fft_f64_...npy
p0_m384_fft_f64_...json
\end{lstlisting}

The patch identifier is parsed from the filename.  The corresponding JSON file
is used to recover the scattering-transform configuration used by the
component-separation run.

Patches can be selected with \texttt{PATCH\_LIST}.  Accepted examples include:

\begin{lstlisting}
export PATCH_LIST="0"
export PATCH_LIST="0,1,2"
export PATCH_LIST="0-15"
\end{lstlisting}

\section{PBC Strategy}

\subsection{Original patches}

The real Planck patches are not periodic.  Their component-separation metadata
therefore records \texttt{pbc=False}.  The target scattering statistics are
computed using this original non-periodic configuration.

\subsection{Periodic synthesis}

The synthetic running maps are generated with:

\begin{lstlisting}
export SYNTHESIS_PBC="1"
export SYNTHESIS_RUNNING_SHAPE="400"
\end{lstlisting}

Thus, a typical synthesis currently uses:

\begin{center}
\begin{tabular}{lll}
\toprule
Object & Shape & Boundary condition \\
\midrule
Component-separated target & $384 \times 384$ & non-periodic \\
Running synthetic map & $400 \times 400$ & periodic \\
Final training pair & $384 \times 384$ & cropped from larger map \\
\bottomrule
\end{tabular}
\end{center}

The larger periodic map permits periodic augmentation operations without
placing the periodic boundary directly inside the final training image.
Augmentation is performed on the larger synthetic map, after which its central
region is cropped back to the original patch dimensions.  Same-patch nuisance
maps are added only after this final crop.

\section{Augmentation}

The augmentation follows the existing Moment Network implementations:

\begin{itemize}
  \item four rotations;
  \item optional horizontal flip; and
  \item periodic translations using \texttt{torch.roll}.
\end{itemize}

Periodic rolling is appropriate for the larger periodic synthetic map.  The
rolled map is cropped back to the target patch size before nuisance addition.
The real nuisance maps are therefore not periodically rolled.

The current test configuration uses:

\begin{lstlisting}
export N_AUGMENTATIONS="8"
export SHIFT_STEP="16"
\end{lstlisting}

The number of augmentations can later be increased to 64, matching the earlier
Moment Network implementations.

\section{Power-Spectrum Constraint}

Initially, synthesis inherited the component-separation metadata setting
\texttt{st\_compute\_ps=True}.  When the running synthesis size was changed
from $384 \times 384$ to $400 \times 400$, this caused a normalization failure.

The number of default power-spectrum bins depends on map size:

\begin{equation}
  n_{\mathrm{bins}} =
  \mathrm{int}\left(2^{\log_2(\min(\mathrm{shape})) - 4}\right).
\end{equation}

Consequently:

\begin{align}
  384 \times 384 &\longrightarrow n_{\mathrm{bins}} = 24, \\
  400 \times 400 &\longrightarrow n_{\mathrm{bins}} = 25.
\end{align}

The target operator stored power-spectrum normalization references with 24
bins, while the running operator produced 25 bins.  PyTorch therefore raised:

\begin{lstlisting}
RuntimeError: The size of tensor a (25) must match
the size of tensor b (24)
\end{lstlisting}

Two changes were made:

\begin{enumerate}
  \item The running operator is explicitly given the target operator's
        \texttt{n\_bins}, ensuring compatibility if power-spectrum constraints
        are enabled.
  \item Power-spectrum constraints are disabled by default for the current
        synthesis experiment:
\end{enumerate}

\begin{lstlisting}
export SYNTHESIS_COMPUTE_PS="0"
\end{lstlisting}

The synthesis therefore currently matches scattering statistics without adding
the power-spectrum terms to the loss.  Power-spectrum constraints can be
restored later by setting \texttt{SYNTHESIS\_COMPUTE\_PS="1"}.

\section{Current Dataset Generation}

For each selected patch, the script currently:

\begin{enumerate}
  \item loads the component-separated Q and U maps;
  \item loads the corresponding JSON metadata;
  \item loads same-patch Q and U nuisance realizations;
  \item synthesizes Q and U independently on $400 \times 400$ periodic maps;
  \item augments the larger periodic synthetic maps;
  \item center-crops augmented maps back to the original patch size;
  \item adds a randomly selected same-patch nuisance realization; and
  \item saves one compressed NumPy dataset shard for the patch.
\end{enumerate}

The current small-run settings are:

\begin{lstlisting}
export PATCH_LIST="0"
export PATCH_LIMIT="1"
export SYNTHESES_PER_PATCH="10"
export SYNTHESIS_BATCH_SIZE="5"
export SYNTHESIS_MAX_ITER="100"
export SYNTHESIS_PBC="1"
export SYNTHESIS_COMPUTE_PS="0"
export SYNTHESIS_RUNNING_SHAPE="400"
export N_AUGMENTATIONS="8"
\end{lstlisting}

\section{Saved Dataset Keys}

Each compressed \texttt{.npz} shard currently contains:

\begin{center}
\begin{tabular}{ll}
\toprule
Key & Contents \\
\midrule
\texttt{xq} & augmented/cropped clean Q targets \\
\texttt{yq} & contaminated Q inputs \\
\texttt{xu} & augmented/cropped clean U targets \\
\texttt{yu} & contaminated U inputs \\
\texttt{sq} & Q syntheses cropped to final patch size \\
\texttt{su} & U syntheses cropped to final patch size \\
\texttt{rq} & raw Q syntheses at running synthesis size \\
\texttt{ru} & raw U syntheses at running synthesis size \\
\texttt{p} & patch identifier \\
\texttt{stem} & source component-separation run stem \\
\texttt{naug} & number of augmentations \\
\texttt{shift} & periodic-roll shift step \\
\texttt{aug} & augmentation description \\
\bottomrule
\end{tabular}
\end{center}

A separate JSON summary is also written for each patch.

\section{Slurm and Cluster Environment}

The scripts are expected to be located on the cluster at:

\begin{lstlisting}
/pscratch/sd/a/atsouros/STL/Moment Network
\end{lstlisting}

The repository root used for importing \texttt{STL\_main} is:

\begin{lstlisting}
/pscratch/sd/a/atsouros/STL
\end{lstlisting}

The sbatch file uses the NERSC account and GPU configuration already used by
the Planck component-separation scripts:

\begin{lstlisting}
#SBATCH -A mp107d_g
#SBATCH --constraint=gpu&hbm80g
#SBATCH --qos=debug
#SBATCH --gpus-per-task=1
\end{lstlisting}

The log filenames are fixed:

\begin{lstlisting}
#SBATCH --output=mn_dataset.log
#SBATCH --error=mn_dataset.err
\end{lstlisting}

The default system Python versions on the login node are too old:

\begin{lstlisting}
python  --version  # Python 2.7.18
python3 --version  # Python 3.6.15
\end{lstlisting}

The sbatch file therefore loads and activates the existing \texttt{cmbenv}
environment before invoking \texttt{python3}.  Unbound-variable checking is
temporarily disabled during environment activation because the conda activation
scripts reference variables that may initially be unset.

\section{Work Still Required}

The current implementation only creates per-patch dataset shards.  The
following stages are still required:

\begin{enumerate}
  \item create a pooled or lazy dataset loader over all patch shards;
  \item define train and validation splits, including held-out-patch
        validation;
  \item train one shared Q network and one shared U network, or evaluate a
        joint two-channel Q/U network;
  \item run inference on every original observed patch;
  \item save posterior mean and posterior standard-deviation maps as
        full-size \texttt{.npy} files; and
  \item evaluate full-map and boundary-sensitive metrics.
\end{enumerate}

\section{Important Validation Considerations}

Randomly splitting augmented samples can place samples originating from the
same patch in both training and validation sets.  This is useful for monitoring
optimization but does not measure generalization to unseen patches.  A
held-out-patch validation split should therefore be used in addition to any
random sample split.

The effect of periodic synthesis should also be evaluated.  Although the final
training maps are cropped from larger periodic maps, comparisons should be made
between full-map metrics and interior-region metrics to detect remaining
boundary effects.

\end{document}
